% Response to the Journal of Forecasting reviewer (manuscript 7441141),
% submitted with the transfer to Applied Stochastic Models in Business and Industry.
\documentclass[12pt]{article}
\usepackage[utf8]{inputenc}
\usepackage{cmap}
\input{glyphtounicode}
\pdfgentounicode=1
\usepackage[margin=1in]{geometry}
\usepackage{mathptmx}
\usepackage{amsmath}
\usepackage[hidelinks]{hyperref}
\usepackage{xurl}
\usepackage{parskip}
\usepackage{xcolor}

\newenvironment{reviewer}{\par\begingroup\itshape\color{black!70}\leftskip=1.5em}{\par\endgroup}

\begin{document}

\begin{center}
{\large\bfseries Response to the review from the Journal of Forecasting}\\[4pt]
Manuscript 7441141, ``Online Anchor-Robust Forecasting under Economic Regime Shift''\\[2pt]
Quang-Vinh Dang
\end{center}

I thank the reviewer for a careful and constructive report. The manuscript was
declined at the \textit{Journal of Forecasting} and has been revised before this
transfer. Below, the reviewer's comments are in italics, each followed by the change
made. Section, table and page numbers refer to the revised manuscript.

\section*{Main concern: the paper is difficult to follow}

\begin{reviewer}
My main concern, however, is that the paper is unusually difficult to follow.
Despite my familiarity with regime-switching models, I struggled to connect the
economic motivation, the assumptions about the anchor, and the empirical exercise.
The paper introduces several extensions before establishing why the basic method is
needed.
\end{reviewer}

I agree, and I have reorganised the whole paper around this point.
\begin{itemize}
\item The introduction now states one question: if a regime shift moves an observed
driver, can the forecaster be made robust to it in advance, and what does that cost
when no shift occurs?
\item The basic method, with the anchor supplied by the analyst, and its guarantee are
presented first (Section~3), with no mention of learning the anchor.
\item The simulation study (Section~4) and the energy application (Section~5) both use
the basic method.
\item Learning the anchor from data and the eigengap diagnostic were previously mixed
into the method section. They are now a separate extension (Section~6), placed after
the main results and introduced as the answer to a practical question raised by
Section~5: what to do when the analyst does not know which driver carries the shift.
\item Section~5 has a new paragraph, ``The anchor and what it assumes''. It states which
indices are used as the anchor, what would have to be true of them for the method to
help, why the news-based indices are more plausible candidates than the VIX, and what
the stress tests of Section~4 predict if the anchor is only weakly aligned with the
true channel. This connects the assumptions of Proposition~1 to the empirical exercise.
\item The conclusion (Section~7) is shorter. It ends with a practical procedure for
deciding whether to use the method.
\end{itemize}

\begin{reviewer}
A simple forecasting example near the beginning would help explain what changes
across regimes, what stays stable, and why the proposed method should improve
forecasts.
\end{reviewer}

The introduction now contains such an example (paragraph ``An example'', Table~1).
\begin{itemize}
\item A target depends on a stable predictor. A second predictor tracks the target
but is also moved by an observed driver, in the way a news-based indicator might rise
with volatility and also with the volume of geopolitical headlines.
\item In the estimation sample the driver has two regimes; in the forecast period it
moves to a level not seen before.
\item The example states explicitly what changes (the relation between the target and
the second predictor) and what stays stable (the relation with the first predictor).
\item Least squares puts weight on the second predictor, has the lower in-sample error,
and roughly doubles its error in the new regime ($0.206$ to $0.452$).
\item Anchor regression removes that weight and loses almost nothing in the new regime
($0.240$ to $0.242$), at a small in-sample cost.
\item The example also shows what the method needs: if the analyst chose a driver
unrelated to the shift, the penalty would remove signal and buy no protection.
\item The script that produces Table~1 is in the replication files
(\texttt{examples/two\_signal\_example.py}).
\end{itemize}

\begin{reviewer}
The paper needs a much clearer account of its purpose and contribution before its
results can be properly assessed.
\end{reviewer}

The paragraph ``What the paper does'' in the introduction now lists the contribution
section by section. The paragraph that follows states the findings, including the
negative ones:
\begin{itemize}
\item the gain in simulation disappears when the anchor is noisy or misaligned;
\item on the energy data OARF is not better than online gradient descent or an
equal-weight average;
\item the learned anchor fails on a minority of streams.
\end{itemize}
The abstract was rewritten in the same way.

\section*{Technical note: the overlapping target}

\begin{reviewer}
The target on pg.~10 is the log of a trailing five-day average of squared daily
returns. At the stated one-step forecast horizon, four of the five squared returns
are already observed. This creates substantial mechanical persistence, which needs to
be considered when interpreting the forecast results.
\end{reviewer}

The reviewer is right, and I confirmed it in the code. The old target shared four of
its five terms with its own lag, and on Brent its correlation with the previous day's
value was $0.89$. The real-data study has been rebuilt and rerun with a target that
does not overlap (Section~5, ``Data and target'').
\begin{itemize}
\item \emph{Target.} The target is now the log of weekly realised variance,
$y_w=\log\sum_{s\in w}r_s^2$, over the daily returns of the week ending on Friday $w$.
\item \emph{Information set.} Every forecast of $y_w$ uses data up to the end of week
$w-1$ only.
\item \emph{No overlap.} Consecutive weeks share no terms. The correlation of the
target with its own lag falls to $0.53$.
\item \emph{Sample.} There are $847$ weeks, of which $678$ are evaluated.
\item \emph{Forecast pool.} The pool was rebuilt for the weekly frequency: last week's
value, HAR-type $4$- and $13$-week means, RiskMetrics and GARCH variances, the lagged
VIX, and a rolling ridge regression.
\item \emph{Pre-processing.} The drivers are no longer pre-scaled with full-sample
statistics, so all scaling now uses past data only.
\end{itemize}
The earlier daily design is kept in Appendix~B, with Table~11, so the reader can see
how much the overlap flattered every method.

\begin{reviewer}
A simple persistence forecast should be reported alongside the proposed method so
that the reader can see what is gained beyond the target's existing predictability.
\end{reviewer}

Persistence (last week's value) is now reported in Table~9, together with the
equal-weight average of the forecast pool, and both are scored under the same
protocol as the learning methods. The two designs give different answers.
\begin{itemize}
\item \emph{Weekly design.} Every learned combination beats persistence. For most
methods the MSE is about a third lower on Brent and a fifth lower on Henry Hub. The
Model Confidence Set excludes persistence on Brent ($p<0.01$). On Henry Hub
persistence stays in the set ($p=0.17$), and the text says that the gain there is not
significant.
\item \emph{Daily design.} In the old daily design persistence is inside the Model
Confidence Set with $p=1.00$ on both series (Appendix~B). This confirms the
reviewer's concern that the earlier results mostly reflected the mechanical
persistence of the target.
\end{itemize}

\begin{reviewer}
A check using a less overlapping target or a longer horizon would also help
establish whether the findings depend on this construction.
\end{reviewer}

The weekly target has no overlap at all. Under it the conclusion about OARF is the
same as before: it is within $1.5\%$ of online gradient descent on both series, and
the Diebold--Mariano test of OARF against online gradient descent gives $p=0.47$ on
Brent and $p=0.54$ on Henry Hub. What changed is the interpretation of the comparison
with simple benchmarks, which now shows a real gain from combining forecasts. The
paper reads the absence of an OARF advantage as consistent with the uncertainty
indices being weak or noisy proxies for the channel, which is the case where the
stress tests of Section~4 predict no gain. It also states that the data cannot rule
out that energy-volatility regimes are not anchor-mediated mean shifts at all.

\section*{Other changes}
\begin{itemize}
\item The appendices on the volatility-timing evaluation (Appendix~C) and on the two
event windows (Appendix~F) were rerun on the weekly design.
\item The text was edited throughout for length and clarity.
\item All code and results are public and regenerate every number in the paper.
\end{itemize}

\end{document}
